\documentclass[11pt]{article}
\usepackage[T1]{fontenc}\usepackage[utf8]{inputenc}\usepackage{lmodern}
\usepackage[spanish,es-noshorthands]{babel}
\usepackage[a4paper,margin=2cm]{geometry}
\usepackage{booktabs,tabularx,xcolor,amssymb}
\usepackage[most]{tcolorbox}
\definecolor{USACHNavy}{RGB}{15,32,67}
\newtcolorbox{guion}{colback=USACHNavy!5,colframe=USACHNavy!60,boxrule=0pt,leftrule=3pt,sharp corners,left=6pt,right=6pt,top=4pt,bottom=4pt}
\setlength{\parindent}{0pt}\setlength{\parskip}{4pt}
\pagestyle{plain}
\begin{document}
{\LARGE\bfseries Guion de Javier: láminas 3, 9 y 10\par}\medskip

\textbf{Tiempo total:} 1 min 45 s. Las frases en negrita son las que el profe dijo en clase que quiere escuchar: díganlas tal cual.\par

\medskip\hrule\medskip

\section*{Lámina 3 $\cdot$ Fuentes de datos (40 s)}

\textit{Entras después de la lámina 2.}\par

\begin{guion}Todo viene del portal del DEMRE, admisión 2026, en cuatro archivos públicos. No trabajamos con una muestra: \textbf{es la población completa}.\par \textbf{La unidad de análisis es la persona}: cada estudiante tiene un ID\_aux, anonimizado, y con esa llave unimos los cuatro archivos.\par De los archivos de inscritos y de puntajes sacamos \textbf{las variables asociadas a la persona}: notas, ranking, puntajes, sexo, ingreso. Y como cada estudiante viene de un colegio, con el \textbf{RBD} construimos \textbf{variables asociadas al colegio}.\par El archivo de postulaciones tiene 1,7 millones de filas porque cada alumno postula a varias carreras; lo dejamos en una fila por estudiante. Con eso \textbf{llegamos al N}, que les muestra mi compañero.\end{guion}

\textit{Pasas la palabra para la lámina 4.}\par

\medskip\hrule\medskip

\section*{Lámina 9 $\cdot$ Descriptivas (25 s)}

\textit{Entras después de la lámina 8.}\par

\begin{guion}Los descriptivos ya \textbf{son coherentes con lo que esperamos a priori}.\par La mediana del promedio de notas es 6,23: los que postulan son estudiantes de buen rendimiento, y por eso el puntaje debería pesar en quedar arriba.\par Y hay 55 \% de mujeres: como ellas tienen mejor NEM pero los hombres más puntaje en matemática, \textbf{el signo del sexo no es obvio}, y es algo que el modelo tiene que resolver.\end{guion}

\medskip\hrule\medskip

\section*{Lámina 10 $\cdot$ Conclusiones (40 s)}

\begin{guion}En resumen: \textbf{llegamos al N de 108.251}, que es la población completa de egresados 2025 que postularon por la vía regular. La Y queda equilibrada, 56 contra 44, porque incluimos a los no matriculados.\par Para la Etapa 2 \textbf{no nos vamos a quedar con las variables disponibles}: vamos a ir al \textbf{MINEDUC y cruzar por RBD} para traer \textbf{matrícula del colegio, número de cursos y ratio estudiantes por profesor}, que nos dicen si el colegio es grande o chico y cuánta atención recibe cada alumno.\par Para medir el entorno preferimos el colegio antes que la comuna, porque \textbf{las comunas son muy heterogéneas}: dentro de una misma comuna hay realidades muy distintas, y el colegio describe mejor el entorno del estudiante.\end{guion}

\textit{Pasas a la lámina 11: «Muchas gracias, ¿preguntas?».}\par

\medskip\hrule\medskip

\section*{Si te preguntan}

\textbf{¿Por qué dicen que es población y no muestra?}\par
\begin{guion}Porque están todos los que cumplen los criterios, no una parte elegida al azar. No hay error muestral.\end{guion}

\textbf{¿Por qué full join?}\par
\begin{guion}Para no perder a nadie al unir: los no matriculados no aparecen en el archivo de matrícula, y los necesitamos para la Y = 0.\end{guion}

\textbf{¿Qué signo esperan de las variables del MINEDUC?}\par
\begin{guion}El ratio estudiantes por profesor, negativo: menos atención por alumno. El tamaño del colegio, ambiguo: más recursos, pero menos atención.\end{guion}
\end{document}
